\chapter{โรคพืชที่เกิดจากจุลินทรีย์และการจัดการ}
\label{chap:ch05}

% ==========================================================
% แผนบริหารการสอนประจำบทที่ 5
% ==========================================================
\section*{แผนบริหารการสอนประจำบทที่ 5}
\addcontentsline{toc}{section}{แผนบริหารการสอนประจำบทที่ 5}

\textbf{หัวข้อเนื้อหาประจำบท}
\begin{enumerate}
    \item มโนทัศน์พื้นฐานของสามเหลี่ยมการเกิดโรคพืช (The Plant Disease Triangle)
    \item วงจรการพัฒนาและการแพร่ระบาดของโรคพืช (Disease Cycle)
    \item โรคพืชที่เกิดจากเชื้อราและโอโอไมซีทในพืชเศรษฐกิจ
    \item โรคพืชที่เกิดจากแบคทีเรียและกลไกการเข้าทำลายเนื้อเยื่อ
    \item การจัดการโรคพืชแบบผสมผสาน (Integrated Disease Management: IDM)
\end{enumerate}

\textbf{วัตถุประสงค์เชิงพฤติกรรม}
\begin{itemize}
    \item อธิบายองค์ประกอบของสามเหลี่ยมการเกิดโรคพืชและความสัมพันธ์เชิงปฏิสัมพันธ์ได้
    \item วินิจฉัยอาการและระบุชนิดของจุลินทรีย์สาเหตุโรคพืชสำคัญในเขตร้อนได้
    \item อธิบายกลไกการเข้าทำลายของเชื้อราสาเหตุโรครากเน่าโคนเน่าและโรคเหี่ยวเขียวได้
    \item ออกแบบแผนการควบคุมและจัดการโรคพืชแบบผสมผสานโดยลดการใช้สารเคมีได้
\end{itemize}

\textbf{กิจกรรมการเรียนการสอน}
\begin{itemize}
    \item การบรรยายและศึกษากล้องจุลทรรศน์ตรวจชิ้นเนื้อเยื่อพืชที่แสดงอาการเป็นโรค
    \item ปฏิบัติการแยกเชื้อบริสุทธิ์จากรอยแผลโรคแอนแทรคโนสและรากเน่าทุเรียน
    \item การศึกษาดูงานแปลงทดสอบการจัดการโรคพืชของคณะเทคโนโลยีการเกษตร
\end{itemize}

\textbf{สื่อการเรียนการสอน}
\begin{itemize}
    \item เอกสารคำสอนบทที่ 5 และภาพแสดงอาการโรคพืชในระดับเซลล์และระดับแปลง
    \item ตัวอย่างจริงของใบมะม่วงที่เป็นแอนแทรคโนส และรากทุเรียนที่เป็นโรครากเน่า
\end{itemize}

\textbf{การประเมินผล}
\begin{itemize}
    \item การประเมินความถูกต้องในการวินิจฉัยและระบุเชื้อสาเหตุโรคจากตัวอย่างพืช
    \item การทดสอบการเขียนแผนจัดการโรคพืชแบบบูรณาการในสวนผลไม้
\end{itemize}

\newpage

% ==========================================================
% เนื้อหาบทเรียน
% ==========================================================

\section{สามเหลี่ยมการเกิดโรคพืชและวงจรการระบาด}
\index{สามเหลี่ยมการเกิดโรคพืช (Plant Disease Triangle)}

การเกิดโรคพืชไม่ได้ขึ้นอยู่กับการมีอยู่ของเชื้อจุลินทรีย์ก่อโรคเพียงอย่างเดียว หากแต่เป็นผลลัพธ์ของปฏิสัมพันธ์ระหว่าง 3 องค์ประกอบหลักที่เกิดขึ้นพร้อมกัน ซึ่งเรียกว่า \textbf{สามเหลี่ยมการเกิดโรคพืช (The Plant Disease Triangle)}:
\begin{enumerate}
    \item \textbf{พืชอาศัยที่อ่อนแอ (Susceptible Host)} พืชพันธุ์ที่ไม่มีความต้านทานทางพันธุกรรม หรือพืชที่อ่อนแอจากสภาวะขาดน้ำ ขาดสารอาหาร หรือมีแผลเปิด
    \item \textbf{เชื้อสาเหตุโรคที่มีความรุนแรง (Virulent Pathogen)} จุลินทรีย์ก่อโรคที่มีปริมาณหัวเชื้อ (Inoculum) เพียงพอและมีความสามารถในการเข้าทำลาย
    \item \textbf{สภาพแวดล้อมที่เหมาะสม (Favorable Environment)} อุณหภูมิ ความชื้นสัมพัทธ์ในอากาศ ความชื้นในดิน และค่า pH ที่เอื้อต่อการงอกของสปอร์และการแพร่ระบาด
\end{enumerate}
หากขาดปัจจัยใดปัจจัยหนึ่งไป โรคพืชจะไม่สามารถเกิดขึ้นหรือลุกลามได้

\begin{figure}[H]
\centering
\begin{tikzpicture}[scale=1.1]
    % Equilateral Triangle
    \coordinate (A) at (0, 3.8);    % Pathogen
    \coordinate (B) at (-3.2, -1.0); % Host
    \coordinate (C) at (3.2, -1.0);  % Environment

    % Triangle fill and outline
    \fill[agriRed!10, draw=agriRed, line width=2pt] (A) -- (B) -- (C) -- cycle;

    % Center circle: Disease occurrence
    \node at (0, 0.5) [circle, draw=agriRed!80!black, fill=white, line width=1.5pt, minimum size=2.2cm, align=center, font=\small\bfseries\color{agriRed!80!black}]
        {การเกิดโรคพืช\\(Disease\\Occurrence)};

    % Vertex A: Pathogen
    \fill[agriRed] (A) circle (6pt);
    \node at (A) [above=8pt, align=center, font=\footnotesize\bfseries\color{agriRed!90!black}]
        {เชื้อสาเหตุโรคที่มีความรุนแรง\\(Virulent Pathogen)};

    % Vertex B: Host
    \fill[agriGreen] (B) circle (6pt);
    \node at (B) [below left=6pt and 0pt, align=center, font=\footnotesize\bfseries\color{agriGreen}]
        {พืชอาศัยที่อ่อนแอ\\(Susceptible Host)};

    % Vertex C: Environment
    \fill[agriBlue] (C) circle (6pt);
    \node at (C) [below right=6pt and 0pt, align=center, font=\footnotesize\bfseries\color{agriBlue}]
        {สภาพแวดล้อมที่เอื้ออำนวย\\(Favorable Environment)};

    % Interacting arrows along edges
    \draw[{Stealth}-{Stealth}, line width=1.2pt, color=agriSlate] (-1.8, 1.6) -- (-0.8, 2.6)
        node[midway, above left, font=\tiny] {การเข้าทำลาย};
    \draw[{Stealth}-{Stealth}, line width=1.2pt, color=agriSlate] (1.8, 1.6) -- (0.8, 2.6)
        node[midway, above right, font=\tiny] {การแพร่กระจาย};
    \draw[{Stealth}-{Stealth}, line width=1.2pt, color=agriSlate] (-1.5, -1.3) -- (1.5, -1.3)
        node[midway, below, font=\tiny] {ความเครียดจากสิ่งแวดล้อม};
\end{tikzpicture}
\caption{สามเหลี่ยมการเกิดโรคพืช (The Plant Disease Triangle) แสดงปฏิสัมพันธ์ขององค์ประกอบสำคัญ}
\label{fig:disease_triangle}
\end{figure}

\section{โรคพืชสำคัญที่เกิดจากเชื้อราและโอโอไมซีท}
\index{โรคพืชจากเชื้อรา (Fungal Plant Diseases)}

เชื้อราและสิ่งมีชีวิตกลุ่มโอโอไมซีท (Oomycetes) เป็นกลุ่มสาเหตุโรคพืชที่มีจำนวนชนิดและก่อความเสียหายทางเศรษฐกิจมากที่สุดในเขตร้อน:

\subsection{1. โรครากเน่าโคนเน่าและผลเน่าทุเรียน (\textit{Phytophthora palmivora})}
\index{ไฟทอปธอรา (Phytophthora palmivora)}
เชื้อ \textit{Phytophthora palmivora} จัดอยู่ในกลุ่มโอโอไมซีท สร้างสปอร์ว่ายน้ำมีแฟลกเจลลา เรียกว่า \textbf{ซูโอสปอร์ (Zoospores)} ซึ่งสามารถว่ายน้ำในฟิล์มน้ำในดินเข้าหาปลายรากพืชตามแรงดึงดูดทางเคมีจากสารหลั่งราก เข้าทำลายเนื้อเยื่อเปลือกและท่อน้ำท่ออาหาร ทำให้โคนต้นมียางไหลสีน้ำตาลแดง เปลือกเน่าเป็นสีดำ ใบเหลืองร่วง และยืนต้นแห้งตายในที่สุด ในช่วงฤดูฝนที่มีความชื้นสัมพัทธ์สูงกว่า 90\% เชื้อสามารถระบาดขึ้นสู่ผลทำให้เกิดโรคผลเน่ารุนแรง

\subsection{2. โรคแอนแทรคโนส (\textit{Colletotrichum} spp.)}
\index{โรคแอนแทรคโนส (Anthracnose)}
เชื้อรา \textit{Colletotrichum gloeosporioides} เข้าทำลายพืชได้หลายชนิด เช่น มะม่วง ทุเรียน พริก และกล้วยไม้ ลักษณะอาการเด่นคือเกิดแผลยุบตัวรูปวงกลมหรือวงรีสีน้ำตาลเข้มถึงดำ บนแผลมีหยดสปอร์สีส้มอมชมพู (Conidial masses) เรียงตัวเป็นวงซ้อนกัน เชื้อสามารถแฝงตัวอยู่ในเนื้อเยื่อพืชแบบสงบนิ่ง (Latent infection) และจะแสดงอาการรุนแรงเมื่อผลไม้เริ่มสุก

\section{โรคพืชที่เกิดจากแบคทีเรีย}
\index{โรคพืชจากแบคทีเรีย (Bacterial Plant Diseases)}

แบคทีเรียก่อโรคพืชไม่มีอวัยวะในการเจาะทะลุผิวเซลล์พืชโดยตรง แต่จะเข้าสู่ต้นพืชผ่านทางบาดแผลธรรมชาติ (รอยหัก รอยรอยแหว่งจากการตัดแต่งกิ่ง) หรือผ่านช่องเปิดธรรมชาติ เช่น ช่องปากใบ (Stomata), เลนทิเซล (Lenticels) และช่องเปิดหยดน้ำ (Hydathodes)
\begin{itemize}
    \item \textbf{โรคเหี่ยวเขียว (\textit{Ralstonia solanacearum})} แบคทีเรียเข้าทำลายท่อน้ำ (Xylem vessels) ของพืชตระกูลมะเขือ พริก และมันฝรั่ง และสร้างเมือกพอลิแซ็กคาไรด์ (Extracellular Polysaccharides: EPS) อุดตันท่อลำเลียงน้ำ ทำให้พืชเหี่ยวเฉาอย่างรวดเร็วทั้งที่ใบยังคงมีสีเขียวสด
    \item \textbf{โรคแคงเกอร์ส้ม (\textit{Xanthomonas citri})} ทำให้เกิดแผลตกสะเก็ดนูนสีน้ำตาล มีวงสีเหลืองล้อมรอบ (Yellow halo) บนใบ กิ่ง และผลส้ม มะนาว
\end{itemize}

\begin{warningbox}{การจัดการโรครากเน่าโคนเน่าทุเรียนแบบผสมผสาน}
ห้ามใช้สารเคมีกำจัดเชื้อราฟอสฟอรัสแอซิดเพียงอย่างเดียว เนื่องจากเชื้อจะพัฒนาความต้านทานได้อย่างรวดเร็ว แนวทางที่ถูกต้องคือ:
1. ปรับปรุงการระบายน้ำในแปลง ไม่ให้น้ำท่วมขังบริเวณโคนต้น\\
2. ปรับค่า pH ของดินให้อยู่ในช่วง 6.0--6.5 ด้วยปูนโดโลไมต์ เพื่อยับยั้งการสร้างซูโอสปอร์\\
3. ทาแผลที่โคนต้นและราดดินด้วยเชื้อราไตรโคเดอร์มาชนิดสดร่วมกับปุ๋ยหมักอินทรีย์อย่างสม่ำเสมอ
\end{warningbox}

\section*{คำถามทบทวนท้ายบทที่ 5}
\addcontentsline{toc}{section}{คำถามทบทวนท้ายบทที่ 5}
\begin{enumerate}
    \item หากต้องการควบคุมการระบาดของโรคพืชโดยการปรับเปลี่ยนสภาพแวดล้อม สามารถทำได้อย่างไรตามหลักการของสามเหลี่ยมการเกิดโรคพืช
    \item จงอธิบายวงจรชีวิตและการแพร่ระบาดของเชื้อ \textit{Phytophthora palmivora} ในสวนทุเรียนช่วงฤดูฝน
    \item การทดสอบการไหลของเมือกแบคทีเรีย (Bacterial Streaming Test) ในน้ำสะอาด ใช้ตรวจวินิจฉัยโรคเหี่ยวเขียวได้อย่างไร
\end{enumerate}

\clearpage